% Chapter 15 infographic -- use right after \label{ch:15}: % Chapter 14 infographic -- use right after \label{ch:14}: % Chapter 13 infographic -- use right after \label{ch:13}: % Chapter 13 infographic -- use right after \label{ch:13}: \input{chapters/ch13-infographic}
\ifcsname FEinfoLoaded\endcsname\else\input{common/fe-infographic}\fi
\begin{FEinfographic}{ch:13}{investors earn the returns their stocks earn, whatever rule they use to decide when to sell}

\infopanel{1}{Three ways a selling rule matters}{%
\infocards
  {The experience}{Even in a fair game, a rule reshapes win rates, tails and the losers still held.}
  {Predictable returns}{With momentum, selling winners and keeping losers collects the wrong returns.}
  {Taxes}{Selling winners pays tax early; selling losers first earns the tax saving.}}

\infopanel{2}{Evidence scoreboard}{%
\begin{infoscore}
\scoreitem{y}{Yes}{Expected return when prices are a fair game}
\scoreitem{n}{No}{The investor's experience}
\scoreitem{n}{No}{Individual investors who sell winners and keep losers}
\scoreitem{n}{No}{After taxes}
\scoreitem{p}{Disputed}{Whole markets, through the timing of money flows}
\scoreitem{y}{Yes}{Patient professional investors}
\end{infoscore}}

\infopanel{3}{The disposition effect in numbers}{%
\infonumbers
  {14.8\%}{of available gains realized (Odean 1998)}
  {9.8\%}{of available losses realized}
  {3.4 pts}{next-year return of winners sold over losers kept}
  {1.3 pts}{a year: NYSE/AMEX dollar-weighted gap (Dichev 2007)}}

\infopanel{4}{Lab: six selling rules, one set of prices}{%
\begin{infolab}
\item \textbf{A.} Optional stopping in simulation. \textbf{B.} Comparing returns over random horizons.
\item \textbf{C.} The rules on a century of 49 industry portfolios.
\item \textbf{D.} Winners sold against losers held. \textbf{E.} Taxes and costs over 30 years.
\end{infolab}}

\infoverdict
  {A selling rule cannot change the expected return of a fair game, but real investors' rules change the experience, collect the wrong returns when returns are predictable, and pay taxes earlier than they need to.}
  {Set the selling rule before buying, ignore the purchase price, realize losses before gains, and compare returns per year invested.}
\end{FEinfographic}

\ifcsname FEinfoLoaded\endcsname\else% ======================================================================
% fe-infographic.tex -- shared layout for the one-page chapter infographics
% Loaded automatically by each chapters/chNN-infographic.tex.
% Works with the Glory style and the agentic style; colours are provided
% here if the preamble does not define them.
% ======================================================================
\providecolor{FEMaroon}{RGB}{128,30,45}
\providecolor{FECream}{RGB}{253,248,238}
\providecolor{FEGold}{RGB}{196,150,60}
\providecolor{fenavy}{HTML}{1B2A6B}
\providecolor{fegray}{HTML}{4D4D4D}
\providecolor{igRisk}{HTML}{1F4E79}
\providecolor{igLuck}{HTML}{8A8A8A}
\providecolor{igSkill}{HTML}{2E6B3F}
\providecolor{igYes}{HTML}{2E6B3F}
\providecolor{igNo}{HTML}{A33A3A}
\providecolor{igPart}{HTML}{C47F00}
\providecolor{igOpen}{HTML}{5D6D7E}

\def\FEinfoLoaded{1}

% Title band and page frame -----------------------------------------
\newenvironment{FEinfographic}[2]{% #1 = chapter label, #2 = claim
  \clearpage\thispagestyle{plain}%
  \begingroup\sffamily\footnotesize\setlength{\parskip}{0pt}\setlength{\parindent}{0pt}%
  \linespread{1.0}\selectfont
  \begin{tcolorbox}[enhanced, colback=FEMaroon, colframe=FEMaroon, arc=3pt,
      boxrule=0pt, left=8pt, right=8pt, top=5pt, bottom=5pt, before skip=0pt, after skip=5pt]
    {\color{white}\bfseries\large Chapter~\ref{#1} at a glance\par}
    \vspace{2pt}{\color{white}\small Claim on trial: \emph{#2}\par}
  \end{tcolorbox}}
  {\par\endgroup\clearpage}

% A numbered panel ----------------------------------------------------
\newcommand{\infopanel}[3]{% #1 number, #2 heading, #3 body
  \begin{tcolorbox}[enhanced, colback=FECream, colframe=FEGold, boxrule=0.6pt, arc=3pt,
      left=7pt, right=7pt, top=4pt, bottom=5pt, before skip=0pt, after skip=5pt]
    {\color{FEMaroon}\bfseries\small #1\quad #2\par}\vspace{3pt}%
    #3
  \end{tcolorbox}}

% Three idea cards ------------------------------------------------------
\newcommand{\infocard}[3]{% #1 colour, #2 title, #3 text
  \begin{minipage}[t]{0.315\linewidth}
    {\setlength{\fboxsep}{2.5pt}\colorbox{#1}{\makebox[\dimexpr\linewidth-2\fboxsep]{\color{white}\bfseries\rule[-0.3ex]{0pt}{2.2ex}#2}}}\par\vspace{3pt}
    {\scriptsize #3\par}
  \end{minipage}}
\newcommand{\infocards}[6]{%
  \infocard{igRisk}{#1}{#2}\hfill\infocard{igLuck}{#3}{#4}\hfill\infocard{igSkill}{#5}{#6}\par}

% Evidence scoreboard ---------------------------------------------------
\newcommand{\vchip}[2]{% #1 = y/n/p/o, #2 = label
  \edef\FEchipc{\ifx y#1igYes\else\ifx n#1igNo\else\ifx p#1igPart\else igOpen\fi\fi\fi}%
  \tcbox[on line, arc=2pt, boxrule=0pt, left=2pt, right=2pt, top=1pt, bottom=1pt,
    colback=\FEchipc,
    fontupper=\sffamily\bfseries\tiny\color{white}]{\makebox[3.6em]{#2}}}
% Items are laid out two to a row.
\newcount\FEscorecount
\newenvironment{infoscore}{\global\FEscorecount=0 \noindent\ignorespaces}{\par}
\newcommand{\scoreitem}[3]{% #1 class, #2 chip label, #3 situation
  \global\advance\FEscorecount by 1
  \begin{minipage}[t]{0.485\linewidth}\raggedright
    \vchip{#1}{#2}\hspace{4pt}\parbox[t]{\dimexpr\linewidth-5.4em}{\scriptsize #3}\vspace{3pt}
  \end{minipage}%
  \ifodd\FEscorecount\hfill\else\par\noindent\fi\ignorespaces}

% Four key numbers ------------------------------------------------------
\newcount\FEnumgroup
\newcommand{\infonum}[2]{%
  \begin{minipage}[t]{0.235\linewidth}
  \begin{tcolorbox}[enhanced, colback=white, colframe=FEGold, boxrule=0.5pt, arc=2pt,
      left=3pt, right=3pt, top=3pt, bottom=3pt, before skip=0pt, after skip=0pt,
      equal height group=FEnum\the\FEnumgroup, valign=top, halign=center]
    {\color{FEMaroon}\bfseries\Large #1\par}\vspace{2pt}{\scriptsize #2\par}
  \end{tcolorbox}
  \end{minipage}}
\newcommand{\infonumbers}[8]{%
  \global\advance\FEnumgroup by 1
  \infonum{#1}{#2}\hfill\infonum{#3}{#4}\hfill\infonum{#5}{#6}\hfill\infonum{#7}{#8}\par}

% Lab parts -------------------------------------------------------------
\newenvironment{infolab}{\scriptsize\begin{itemize}[leftmargin=1.1em, itemsep=1pt, topsep=1pt, parsep=0pt]}{\end{itemize}}

% Verdict band ----------------------------------------------------------
\newcommand{\infoverdict}[2]{% #1 = where the evidence stands, #2 = what investors can use
  \begin{tcolorbox}[enhanced, colback=white, colframe=FEMaroon, boxrule=0.8pt, arc=3pt,
      leftrule=4pt, left=7pt, right=7pt, top=4pt, bottom=5pt, before skip=0pt, after skip=0pt]
    {\color{FEMaroon}\bfseries\small Where the evidence stands\par}\vspace{2pt}
    {\scriptsize #1\par}\vspace{4pt}
    {\color{FEMaroon}\bfseries\small What investors can use\par}\vspace{2pt}
    {\scriptsize #2\par}
  \end{tcolorbox}}
\fi
\begin{FEinfographic}{ch:13}{investors earn the returns their stocks earn, whatever rule they use to decide when to sell}

\infopanel{1}{Three ways a selling rule matters}{%
\infocards
  {The experience}{Even in a fair game, a rule reshapes win rates, tails and the losers still held.}
  {Predictable returns}{With momentum, selling winners and keeping losers collects the wrong returns.}
  {Taxes}{Selling winners pays tax early; selling losers first earns the tax saving.}}

\infopanel{2}{Evidence scoreboard}{%
\begin{infoscore}
\scoreitem{y}{Yes}{Expected return when prices are a fair game}
\scoreitem{n}{No}{The investor's experience}
\scoreitem{n}{No}{Individual investors who sell winners and keep losers}
\scoreitem{n}{No}{After taxes}
\scoreitem{p}{Disputed}{Whole markets, through the timing of money flows}
\scoreitem{y}{Yes}{Patient professional investors}
\end{infoscore}}

\infopanel{3}{The disposition effect in numbers}{%
\infonumbers
  {14.8\%}{of available gains realized (Odean 1998)}
  {9.8\%}{of available losses realized}
  {3.4 pts}{next-year return of winners sold over losers kept}
  {1.3 pts}{a year: NYSE/AMEX dollar-weighted gap (Dichev 2007)}}

\infopanel{4}{Lab: six selling rules, one set of prices}{%
\begin{infolab}
\item \textbf{A.} Optional stopping in simulation. \textbf{B.} Comparing returns over random horizons.
\item \textbf{C.} The rules on a century of 49 industry portfolios.
\item \textbf{D.} Winners sold against losers held. \textbf{E.} Taxes and costs over 30 years.
\end{infolab}}

\infoverdict
  {A selling rule cannot change the expected return of a fair game, but real investors' rules change the experience, collect the wrong returns when returns are predictable, and pay taxes earlier than they need to.}
  {Set the selling rule before buying, ignore the purchase price, realize losses before gains, and compare returns per year invested.}
\end{FEinfographic}

\ifcsname FEinfoLoaded\endcsname\else% ======================================================================
% fe-infographic.tex -- shared layout for the one-page chapter infographics
% Loaded automatically by each chapters/chNN-infographic.tex.
% Works with the Glory style and the agentic style; colours are provided
% here if the preamble does not define them.
% ======================================================================
\providecolor{FEMaroon}{RGB}{128,30,45}
\providecolor{FECream}{RGB}{253,248,238}
\providecolor{FEGold}{RGB}{196,150,60}
\providecolor{fenavy}{HTML}{1B2A6B}
\providecolor{fegray}{HTML}{4D4D4D}
\providecolor{igRisk}{HTML}{1F4E79}
\providecolor{igLuck}{HTML}{8A8A8A}
\providecolor{igSkill}{HTML}{2E6B3F}
\providecolor{igYes}{HTML}{2E6B3F}
\providecolor{igNo}{HTML}{A33A3A}
\providecolor{igPart}{HTML}{C47F00}
\providecolor{igOpen}{HTML}{5D6D7E}

\def\FEinfoLoaded{1}

% Title band and page frame -----------------------------------------
\newenvironment{FEinfographic}[2]{% #1 = chapter label, #2 = claim
  \clearpage\thispagestyle{plain}%
  \begingroup\sffamily\footnotesize\setlength{\parskip}{0pt}\setlength{\parindent}{0pt}%
  \linespread{1.0}\selectfont
  \begin{tcolorbox}[enhanced, colback=FEMaroon, colframe=FEMaroon, arc=3pt,
      boxrule=0pt, left=8pt, right=8pt, top=5pt, bottom=5pt, before skip=0pt, after skip=5pt]
    {\color{white}\bfseries\large Chapter~\ref{#1} at a glance\par}
    \vspace{2pt}{\color{white}\small Claim on trial: \emph{#2}\par}
  \end{tcolorbox}}
  {\par\endgroup\clearpage}

% A numbered panel ----------------------------------------------------
\newcommand{\infopanel}[3]{% #1 number, #2 heading, #3 body
  \begin{tcolorbox}[enhanced, colback=FECream, colframe=FEGold, boxrule=0.6pt, arc=3pt,
      left=7pt, right=7pt, top=4pt, bottom=5pt, before skip=0pt, after skip=5pt]
    {\color{FEMaroon}\bfseries\small #1\quad #2\par}\vspace{3pt}%
    #3
  \end{tcolorbox}}

% Three idea cards ------------------------------------------------------
\newcommand{\infocard}[3]{% #1 colour, #2 title, #3 text
  \begin{minipage}[t]{0.315\linewidth}
    {\setlength{\fboxsep}{2.5pt}\colorbox{#1}{\makebox[\dimexpr\linewidth-2\fboxsep]{\color{white}\bfseries\rule[-0.3ex]{0pt}{2.2ex}#2}}}\par\vspace{3pt}
    {\scriptsize #3\par}
  \end{minipage}}
\newcommand{\infocards}[6]{%
  \infocard{igRisk}{#1}{#2}\hfill\infocard{igLuck}{#3}{#4}\hfill\infocard{igSkill}{#5}{#6}\par}

% Evidence scoreboard ---------------------------------------------------
\newcommand{\vchip}[2]{% #1 = y/n/p/o, #2 = label
  \edef\FEchipc{\ifx y#1igYes\else\ifx n#1igNo\else\ifx p#1igPart\else igOpen\fi\fi\fi}%
  \tcbox[on line, arc=2pt, boxrule=0pt, left=2pt, right=2pt, top=1pt, bottom=1pt,
    colback=\FEchipc,
    fontupper=\sffamily\bfseries\tiny\color{white}]{\makebox[3.6em]{#2}}}
% Items are laid out two to a row.
\newcount\FEscorecount
\newenvironment{infoscore}{\global\FEscorecount=0 \noindent\ignorespaces}{\par}
\newcommand{\scoreitem}[3]{% #1 class, #2 chip label, #3 situation
  \global\advance\FEscorecount by 1
  \begin{minipage}[t]{0.485\linewidth}\raggedright
    \vchip{#1}{#2}\hspace{4pt}\parbox[t]{\dimexpr\linewidth-5.4em}{\scriptsize #3}\vspace{3pt}
  \end{minipage}%
  \ifodd\FEscorecount\hfill\else\par\noindent\fi\ignorespaces}

% Four key numbers ------------------------------------------------------
\newcount\FEnumgroup
\newcommand{\infonum}[2]{%
  \begin{minipage}[t]{0.235\linewidth}
  \begin{tcolorbox}[enhanced, colback=white, colframe=FEGold, boxrule=0.5pt, arc=2pt,
      left=3pt, right=3pt, top=3pt, bottom=3pt, before skip=0pt, after skip=0pt,
      equal height group=FEnum\the\FEnumgroup, valign=top, halign=center]
    {\color{FEMaroon}\bfseries\Large #1\par}\vspace{2pt}{\scriptsize #2\par}
  \end{tcolorbox}
  \end{minipage}}
\newcommand{\infonumbers}[8]{%
  \global\advance\FEnumgroup by 1
  \infonum{#1}{#2}\hfill\infonum{#3}{#4}\hfill\infonum{#5}{#6}\hfill\infonum{#7}{#8}\par}

% Lab parts -------------------------------------------------------------
\newenvironment{infolab}{\scriptsize\begin{itemize}[leftmargin=1.1em, itemsep=1pt, topsep=1pt, parsep=0pt]}{\end{itemize}}

% Verdict band ----------------------------------------------------------
\newcommand{\infoverdict}[2]{% #1 = where the evidence stands, #2 = what investors can use
  \begin{tcolorbox}[enhanced, colback=white, colframe=FEMaroon, boxrule=0.8pt, arc=3pt,
      leftrule=4pt, left=7pt, right=7pt, top=4pt, bottom=5pt, before skip=0pt, after skip=0pt]
    {\color{FEMaroon}\bfseries\small Where the evidence stands\par}\vspace{2pt}
    {\scriptsize #1\par}\vspace{4pt}
    {\color{FEMaroon}\bfseries\small What investors can use\par}\vspace{2pt}
    {\scriptsize #2\par}
  \end{tcolorbox}}
\fi
\begin{FEinfographic}{ch:14}{a strategy that is profitable in a backtest is profitable to trade}

\infopanel{1}{Three costs a backtest leaves out}{%
\infocards
  {Turnover}{Every trade pays half the bid--ask spread, and fast strategies trade most of their portfolio each month.}
  {Market impact}{Large trades move prices against the trader, and the cost grows with the size of the fund.}
  {The short side}{Borrowing fees, recalls and bans make many anomalies costly to short.}}

\infopanel{2}{Evidence scoreboard}{%
\begin{infoscore}
\scoreitem{y}{Yes}{Slow strategies, such as value}
\scoreitem{p}{Partly}{Momentum}
\scoreitem{n}{No}{Short-term reversal}
\scoreitem{n}{No}{High-turnover strategies in general}
\scoreitem{y}{More often}{For a large, patient institution}
\scoreitem{n}{Often not}{Strategies that rely on the short side}
\scoreitem{n}{No}{The average anomaly}
\end{infoscore}}

\infopanel{3}{Monthly turnover and what it costs}{%
\infonumbers
  {2\%}{monthly turnover of value}
  {25\%}{monthly turnover of momentum}
  {76\%}{monthly turnover of short-term reversal}
  {4 bp}{a month: the average anomaly after costs and decay}}

\infopanel{4}{Lab: replicating Novy-Marx and Velikov (2016)}{%
\begin{infolab}
\item \textbf{A.} Gross returns of value, momentum and short-term reversal. \textbf{B.} Turnover in a simulated universe.
\item \textbf{C.} Net returns and break-even costs from 0 to 100 basis points.
\item \textbf{D.} Capacity as assets grow to \$1 trillion. \textbf{E.} Buy/hold bands that cut momentum's turnover in half.
\end{infolab}}

\infoverdict
  {Whether a profitable backtest is profitable to trade depends on how much the strategy trades, what it trades and how much money is behind it: slow strategies in liquid stocks survive realistic costs, while fast strategies in small stocks often do not.}
  {Ask for turnover, the break-even cost against what you will pay, the cost of the short side and the net return at your planned size.}
\end{FEinfographic}

\ifcsname FEinfoLoaded\endcsname\else% ======================================================================
% fe-infographic.tex -- shared layout for the one-page chapter infographics
% Loaded automatically by each chapters/chNN-infographic.tex.
% Works with the Glory style and the agentic style; colours are provided
% here if the preamble does not define them.
% ======================================================================
\providecolor{FEMaroon}{RGB}{128,30,45}
\providecolor{FECream}{RGB}{253,248,238}
\providecolor{FEGold}{RGB}{196,150,60}
\providecolor{fenavy}{HTML}{1B2A6B}
\providecolor{fegray}{HTML}{4D4D4D}
\providecolor{igRisk}{HTML}{1F4E79}
\providecolor{igLuck}{HTML}{8A8A8A}
\providecolor{igSkill}{HTML}{2E6B3F}
\providecolor{igYes}{HTML}{2E6B3F}
\providecolor{igNo}{HTML}{A33A3A}
\providecolor{igPart}{HTML}{C47F00}
\providecolor{igOpen}{HTML}{5D6D7E}

\def\FEinfoLoaded{1}

% Title band and page frame -----------------------------------------
\newenvironment{FEinfographic}[2]{% #1 = chapter label, #2 = claim
  \clearpage\thispagestyle{plain}%
  \begingroup\sffamily\footnotesize\setlength{\parskip}{0pt}\setlength{\parindent}{0pt}%
  \linespread{1.0}\selectfont
  \begin{tcolorbox}[enhanced, colback=FEMaroon, colframe=FEMaroon, arc=3pt,
      boxrule=0pt, left=8pt, right=8pt, top=5pt, bottom=5pt, before skip=0pt, after skip=5pt]
    {\color{white}\bfseries\large Chapter~\ref{#1} at a glance\par}
    \vspace{2pt}{\color{white}\small Claim on trial: \emph{#2}\par}
  \end{tcolorbox}}
  {\par\endgroup\clearpage}

% A numbered panel ----------------------------------------------------
\newcommand{\infopanel}[3]{% #1 number, #2 heading, #3 body
  \begin{tcolorbox}[enhanced, colback=FECream, colframe=FEGold, boxrule=0.6pt, arc=3pt,
      left=7pt, right=7pt, top=4pt, bottom=5pt, before skip=0pt, after skip=5pt]
    {\color{FEMaroon}\bfseries\small #1\quad #2\par}\vspace{3pt}%
    #3
  \end{tcolorbox}}

% Three idea cards ------------------------------------------------------
\newcommand{\infocard}[3]{% #1 colour, #2 title, #3 text
  \begin{minipage}[t]{0.315\linewidth}
    {\setlength{\fboxsep}{2.5pt}\colorbox{#1}{\makebox[\dimexpr\linewidth-2\fboxsep]{\color{white}\bfseries\rule[-0.3ex]{0pt}{2.2ex}#2}}}\par\vspace{3pt}
    {\scriptsize #3\par}
  \end{minipage}}
\newcommand{\infocards}[6]{%
  \infocard{igRisk}{#1}{#2}\hfill\infocard{igLuck}{#3}{#4}\hfill\infocard{igSkill}{#5}{#6}\par}

% Evidence scoreboard ---------------------------------------------------
\newcommand{\vchip}[2]{% #1 = y/n/p/o, #2 = label
  \edef\FEchipc{\ifx y#1igYes\else\ifx n#1igNo\else\ifx p#1igPart\else igOpen\fi\fi\fi}%
  \tcbox[on line, arc=2pt, boxrule=0pt, left=2pt, right=2pt, top=1pt, bottom=1pt,
    colback=\FEchipc,
    fontupper=\sffamily\bfseries\tiny\color{white}]{\makebox[3.6em]{#2}}}
% Items are laid out two to a row.
\newcount\FEscorecount
\newenvironment{infoscore}{\global\FEscorecount=0 \noindent\ignorespaces}{\par}
\newcommand{\scoreitem}[3]{% #1 class, #2 chip label, #3 situation
  \global\advance\FEscorecount by 1
  \begin{minipage}[t]{0.485\linewidth}\raggedright
    \vchip{#1}{#2}\hspace{4pt}\parbox[t]{\dimexpr\linewidth-5.4em}{\scriptsize #3}\vspace{3pt}
  \end{minipage}%
  \ifodd\FEscorecount\hfill\else\par\noindent\fi\ignorespaces}

% Four key numbers ------------------------------------------------------
\newcount\FEnumgroup
\newcommand{\infonum}[2]{%
  \begin{minipage}[t]{0.235\linewidth}
  \begin{tcolorbox}[enhanced, colback=white, colframe=FEGold, boxrule=0.5pt, arc=2pt,
      left=3pt, right=3pt, top=3pt, bottom=3pt, before skip=0pt, after skip=0pt,
      equal height group=FEnum\the\FEnumgroup, valign=top, halign=center]
    {\color{FEMaroon}\bfseries\Large #1\par}\vspace{2pt}{\scriptsize #2\par}
  \end{tcolorbox}
  \end{minipage}}
\newcommand{\infonumbers}[8]{%
  \global\advance\FEnumgroup by 1
  \infonum{#1}{#2}\hfill\infonum{#3}{#4}\hfill\infonum{#5}{#6}\hfill\infonum{#7}{#8}\par}

% Lab parts -------------------------------------------------------------
\newenvironment{infolab}{\scriptsize\begin{itemize}[leftmargin=1.1em, itemsep=1pt, topsep=1pt, parsep=0pt]}{\end{itemize}}

% Verdict band ----------------------------------------------------------
\newcommand{\infoverdict}[2]{% #1 = where the evidence stands, #2 = what investors can use
  \begin{tcolorbox}[enhanced, colback=white, colframe=FEMaroon, boxrule=0.8pt, arc=3pt,
      leftrule=4pt, left=7pt, right=7pt, top=4pt, bottom=5pt, before skip=0pt, after skip=0pt]
    {\color{FEMaroon}\bfseries\small Where the evidence stands\par}\vspace{2pt}
    {\scriptsize #1\par}\vspace{4pt}
    {\color{FEMaroon}\bfseries\small What investors can use\par}\vspace{2pt}
    {\scriptsize #2\par}
  \end{tcolorbox}}
\fi
\begin{FEinfographic}{ch:15}{do low-risk stocks earn higher risk-adjusted returns than high-risk stocks?}

\infopanel{1}{The low-risk anomaly}{%
\infocards
  {A flat SML}{Returns rise with beta far less than the CAPM predicts, so low-beta stocks have positive alphas.}
  {Betting against beta}{Lever low-beta stocks and short high-beta ones to build a market-neutral portfolio.}
  {Why?}{Leverage constraints, lottery preferences, benchmarking and costly short selling all have support.}}

\infopanel{2}{Where the evidence stands: what published studies found}{%
\begin{infoscore}
\scoreitem{y}{Yes}{A flat security market line}
\scoreitem{y}{Yes}{Idiosyncratic volatility}
\scoreitem{p}{Partly}{The beta-neutral BAB factor}
\scoreitem{p}{Weakly}{After profitability and investment factors}
\scoreitem{p}{Weakly}{With betas that vary over time}
\scoreitem{p}{Partly}{Interest-rate exposure}
\scoreitem{o}{Disputed}{Why the line is flat}
\end{infoscore}}

\infopanel{3}{Key numbers}{%
\infonumbers
  {about 1\%}{a month: most minus least volatile stocks}
  {\$1.05}{of each BAB dollar in the smallest 1\% of stocks}
  {0.7--0.8}{typical beta of a minimum-volatility fund}
  {73.5\%}{of CFOs who use the CAPM for the cost of equity}}

\infopanel{4}{Lab: The security market line and betting against beta}{%
\begin{infolab}
\item \textbf{A.} The security market line for portfolios sorted on beta, volatility and residual volatility.
\item \textbf{B.} Low-minus-high and BAB-style long--short portfolios.
\item \textbf{C.} Weighting, subperiods, interest rates and crises. \textbf{D.} Low-risk quintiles from an investor's view.
\end{infolab}}

\infosummary
  {The security market line is flatter than the CAPM predicts, one of the most robust findings in empirical finance. Whether low risk earns a separate premium that a beta-neutral portfolio can collect depends on construction, small stocks and overlap with other factors.}
  {Minimum-volatility funds offer a better return per unit of risk, not higher returns, and can lag the market for years; BAB funds are levered, market-neutral and crash-prone.}
\end{FEinfographic}
